\documentclass[letterpaper,journal]{IEEEtran}
\usepackage{amsmath,amsfonts}
\usepackage{algorithmic}
\usepackage{algorithm}
\usepackage{array}
\usepackage[caption=false,font=normalsize,labelfont=sf,textfont=sf]{subfig}
\usepackage{textcomp}
\usepackage{stfloats}
\usepackage{url}
\usepackage{verbatim}
\usepackage{graphicx}
\usepackage{cite}
\hyphenation{op-tical net-works semi-conduc-tor IEEE-Xplore}

\begin{document}

\title{Delay-Robust Closed-Loop Inverse Kinematics Control for Unified 25-DoF Arm-Hand Dexterous Teleoperation Systems}

\author{Thien Bao Tran, The Tri Bui, Huu Tran Nhat Le, and Ha Quang Thinh Ngo
\thanks{Ha Quang Thinh Ngo is with the FPT University, Ho Chi Minh City, Vietnam (corresponding author, phone: +84-961-382007; e-mail: thinhnhq2@fe.edu.vn).}%
\thanks{Thien Bao Tran is with the Ho Chi Minh City University of Technology (HCMUT), 268 Ly Thuong Kiet, Dien Hong Ward, Ho Chi Minh City, Vietnam (e-mail: bao.tran2312sdh@hcmut.edu.vn).}%
\thanks{The Tri Bui is with the Ho Chi Minh City University of Technology (HCMUT), 268 Ly Thuong Kiet, Dien Hong Ward, Ho Chi Minh City, Vietnam (e-mail: bttri.sdh231@hcmut.edu.vn).}%
\thanks{Huu Tran Nhat Le is with the Ho Chi Minh City University of Technology (HCMUT), 268 Ly Thuong Kiet, Dien Hong Ward, Ho Chi Minh City, Vietnam (e-mail: lhtnhat.sdh232@hcmut.edu.vn).}}

\markboth{IEEE Transactions on Robotics,~Vol.~22, No.~7, July~2026}%
{Tran \MakeLowercase{\textit{et al.}}: Delay-Robust Closed-Loop Inverse Kinematics Control for Unified 25-DoF Arm-Hand Systems}

\maketitle

\begin{abstract}
This paper presents a delay-robust Closed-Loop Inverse Kinematics (CLIK) control scheme for a unified 25-DoF arm-hand robotic system (6-DoF UR5 manipulator and 19-DoF custom dexterous hand) under time-varying communication delays. By introducing a filtered reference trajectory and modeling latency under a true-and-noise framework, we prove Local Input-to-State Stability (LISS) using a Lyapunov-Krasovskii Functional (LKF) and Linear Matrix Inequalities (LMIs). A damped least-squares (DLS) Jacobian inverse is utilized to guarantee singularity-robust tracking, while secondary null-space projection optimizes joint limits, self-collision avoidance, and natural anthropomorphic postures. The effectiveness of the proposed method is validated at 100 Hz in a physical simulator.
\end{abstract}

\begin{IEEEkeywords}
Closed-Loop Inverse Kinematics, Lyapunov-Krasovskii Functional, Input-to-State Stability, Teleoperation, Redundant Robots.
\end{IEEEkeywords}

\section{Introduction}
\IEEEPARstart{V}{ision-based} teleoperation systems have emerged as a cornerstone for enabling human operators to perform complex tasks using highly dexterous robotic arm-hand platforms. Recently, the robotics community has seen a paradigm shift toward learning-based visuomotor policies, such as the 3D Diffusion Policy \cite{ref_diffpolicy3d} and 3D Flow Diffusion Policy \cite{ref_flowdiff}, which demand large volumes of high-fidelity, smooth, and collision-free human demonstrations. Standard frameworks like AnyTeleop \cite{ref_anyteleop} and DexSim2Real \cite{ref_dexsim2real} utilize markerless vision-based tracking pipelines, leveraging egocentric or multi-camera hand pose estimation \cite{ref_egocentric3d, ref_deltadorsal, ref_egohand_bench} and advanced deep networks such as A2J-Transformers \cite{ref_a2j_trans} and Deformer \cite{ref_deformer} to track human hand keypoints without wearable gloves. 

Despite these vision advancements, deploying such systems in remote, network-coupled environments introduces significant challenges. Data transmission over wireless channels (Wi-Fi, 4G/5G, or internet connections) inevitably suffers from time-varying communication latencies, packet dropouts, and non-differentiable network jitter \cite{ref_tase_tele, ref_tcyb_nn, ref_tmech_manip, ref_tcst_delay}. From a control theory perspective, packet jitter creates high-frequency non-differentiable delay trajectories, rendering standard Closed-Loop Inverse Kinematics (CLIK) velocity feedforward terms mathematically undefined. Furthermore, the 25-DoF unified arm-hand chain (comprising a 6-DoF UR5 manipulator and a 19-DoF custom multi-fingered hand) is highly redundant and prone to kinematic singularities, joint limit violations, and self-collisions.

To address these limitations, this paper proposes a delay-robust and singularity-robust CLIK control scheme. The primary contributions of this work are three-fold:
\begin{itemize}
\item We formulate a true-and-noise delay model that isolates non-differentiable network jitter into a bounded spatial tracking offset, enabling smooth tracking using a filtered delay reference.
\item We prove local Input-to-State Stability (LISS) of the tracking error under time-varying delays using a Lyapunov-Krasovskii Functional (LKF) and Linear Matrix Inequalities (LMIs) derived via the reciprocal convexity lemma.
\item We integrate a Damped Least-Squares (DLS) Jacobian inverse with secondary null-space projection tasks to systematically handle kinematic singularities, joint limits, self-collisions, and human-like hand postures.
\end{itemize}

\section{Related Work}
\subsection{Dexterous Hand-Arm Retargeting and Control}
Mapping human hand-arm motion to high-DoF robotic platforms requires solving complex joint retargeting and collision avoidance problems \cite{ref_cybernetics_armhand, ref_xin_retarget}. Many research efforts focus on optimization-based retargeting to transfer finger joint angles from human hands to robotic counterparts \cite{ref_tro_grasp, ref_dexhil, ref_zero_shot_dex}. However, integrating 19-DoF hands with a 6-DoF arm adds significant redundancy, requiring null-space projection techniques to maintain natural anthropomorphic posture configurations \cite{ref_amp}. Singularity handling is also critical; standard pseudoinverses suffer from joint velocity explosion near singular configurations, prompting the use of Damped Least-Squares (DLS) algorithms \cite{ref_neuro_imped}.

\subsection{Bilateral Teleoperation and Delayed Robot Control}
Time-delay compensation is a classical problem in bilateral teleoperation systems. Various methodologies have been developed to guarantee closed-loop stability, including sliding mode control \cite{ref_iecon_sliding}, neural network-based adaptive control \cite{ref_tcyb_nn, ref_ojies_coor, ref_tsmc_nn}, predictor-based schemes \cite{ref_access_pred}, and event-triggered communication protocols \cite{ref_tcyb_event, ref_cac_event}. To analyze the stability of delayed systems, Lyapunov-Krasovskii Functional (LKF) techniques combined with LMI constraints are widely used \cite{ref_tcs_haptic, ref_tsmc_linear, ref_tcyb_piece, ref_ijsr_tele, ref_tcst_delayed_force, ref_tro_rcm}. However, most delayed control methods are limited to 6-DoF or 7-DoF manipulators, and do not scale to unified 25-DoF arm-hand systems. 

A State-of-the-Art (SOTA) comparison of the proposed scheme against other representative teleoperation control frameworks is provided in Table \ref{tab:sota}.

\begin{table*}[t]
\centering
\caption{State-of-the-Art (SOTA) Comparison of Robotic Teleoperation Controllers}
\label{tab:sota}
\begin{tabular}{|l|c|c|c|c|c|}
\hline
\textbf{Framework} & \textbf{Target System} & \textbf{DoF} & \textbf{Delay Handling} & \textbf{Stability Analysis} & \textbf{Secondary Constraints} \\
\hline
Li \textit{et al.} \cite{ref_cybernetics_armhand} & Arm-Hand System & High (25 DoF) & No Delay & No Stability Proof & Joint Limits Only \\
AnyTeleop \cite{ref_anyteleop} & Unified Arm-Hand & High (26 DoF) & No Delay & Kinematic Heuristics Only & Joint Limits \& Collision \\
Jing \textit{et al.} \cite{ref_tase_tele} & Bilateral Manipulator & Low (2 DoF) & Time-Varying Delay & Lyapunov Stability & None \\
Bertino \textit{et al.} \cite{ref_tcst_delay} & 7-DoF Manipulator & Low (7 DoF) & Constant Input Delay & Delay-Adaptive Lyapunov & Heuristic Limits \\
\textbf{Proposed Scheme} & \textbf{Unified Arm-Hand} & \textbf{High (25 DoF)} & \textbf{Time-Varying + Jitter} & \textbf{LKF-LMI Rigorous LISS} & \textbf{Limits, Collision \& Synergy} \\
\hline
\end{tabular}
\end{table*}

\section{Problem Formulation}

Consider an integrated robotic system consisting of a 6-DoF manipulator (UR5) and a 19-DoF custom dexterous hand, constituting a highly redundant kinematic chain with a joint configuration vector $q(t) \in \mathbb{R}^n$ ($n = 25$), defined as:
\begin{equation}
q(t) = \begin{bmatrix} q_a(t) \\ q_h(t) \end{bmatrix}
\end{equation}
where $q_a(t) \in \mathbb{R}^6$ represents the joint angles of the robotic arm, and $q_h(t) \in \mathbb{R}^{19}$ denotes the active joint angles of the multi-fingered hand. 

The forward kinematics mapping the joint space configuration to the task space pose $x(t) \in \mathcal{X}$ is given by:
\begin{equation}
x(t) = f(q(t))
\end{equation}
where the task space $\mathcal{X} = \mathbb{R}^3 \times \mathcal{S}^3 \times \mathbb{R}^{n_f}$ is designed with $n_f = 5$ relative Euclidean finger opposition distances (specifically, thumb-index, thumb-middle, thumb-ring, thumb-pinky, and index-pinky opposition pairs) to prevent kinematic over-constraint ($m = 3 + 3 + 5 = 11$ task-space dimensions, which is well below $n = 25$ total joints). 

Differentiating the forward kinematics yields the differential kinematics:
\begin{equation}
v(t) = J(q(t))\dot{q}(t)
\end{equation}
where $v(t) = [\dot{p}_w^T(t), \omega_w^T(t), \dot{d}_f^T(t)]^T \in \mathbb{R}^m$ is the spatial task-space velocity vector ($\omega_w$ being the physical angular velocity of the wrist), and $J(q(t)) \in \mathbb{R}^{m \times n}$ is the analytical system Jacobian.

In a remote teleoperation setup, the desired trajectory $x_d(t)$ is commanded by a human operator through a vision-based hand tracker. The controller receives $x_d(t)$ under an unknown but deterministic time-varying communication delay $\tau(t)$, satisfying the following conditions:

\textbf{Assumption 1}: The communication delay $\tau(t)$ is continuous and bounded such that:
\begin{equation}
0 \le \tau(t) \le \tau_m, \quad \forall t \ge 0
\end{equation}
where $\tau_m$ is a known upper bound of the latency.

\textbf{Assumption 2}: The rate of change of the true delay is bounded on both sides to prevent unstable tracking:
\begin{equation}
|\dot{\tau}_{true}(t)| \le d < 1, \quad \forall t \ge 0
\end{equation}

\textbf{Assumption 3}: The measured latency consists of a true delay $\tau_{true}(t)$ satisfying Assumptions 1 and 2, and a high-frequency bounded measurement noise $n(t)$ representing non-differentiable network packet jitter:
\begin{equation}
\tau(t) = \tau_{true}(t) + n(t), \quad \|n(t)\| \le n_{max}
\end{equation}

Because the raw delayed trajectory $x_d(t - \tau(t))$ contains the high-frequency non-differentiable jitter $n(t)$, directly feeding it into a continuous-time differential controller violates smoothness requirements. In practice, the reference signal is processed through a continuous-time low-pass filter (or jitter buffer). Thus, we define the \textit{filtered reference trajectory} $x_{d,filt}(t)$ equivalent to evaluating the trajectory at the filtered delay $\tau_{filt}(t)$:
\begin{equation}
x_{d,filt}(t) = x_d(t - \tau_{filt}(t))
\end{equation}
where $\tau_{filt}(t)$ is generated such that it is smooth, strictly bounded ($0 \le \tau_{filt}(t) \le \tau_m$), and its derivative is strictly bounded ($|\dot{\tau}_{filt}(t)| \le d < 1$).

Under delayed teleoperation, the task-space tracking error $e(t) \in \mathbb{R}^m$ is redefined relative to the filtered reference:
\begin{equation}
e(t) = \begin{bmatrix} e_p(t) \\ e_o(t) \\ e_f(t) \end{bmatrix}
\end{equation}
where:
\begin{align}
e_p(t) &= p_{d}(t-\tau_{filt}(t)) - p_w(t), \\
e_o(t) &= \eta_d(t-\tau_{filt}(t))\epsilon(t) - \eta(t)\epsilon_d(t-\tau_{filt}(t)) \nonumber \\
&\quad - \epsilon^\times(t)\epsilon_d(t-\tau_{filt}(t)), \\
e_f(t) &= d_{d}(t-\tau_{filt}(t)) - d_f(t),
\end{align}
and $e_o(t) \in \mathbb{R}^3$ represents the quaternion-based orientation error, with $\epsilon^\times$ being the skew-symmetric matrix of $\epsilon$. By redefining the tracking error relative to the filtered reference, the physical jitter $n(t)$ is decoupled from the closed-loop differential kinematics, acting instead as a bounded spatial tracking offset (since $x_d$ is Lipschitz continuous and $n(t)$ is bounded by Assumption 3).

The primary control objective is to design a joint velocity controller $\dot{q}(t)$ that guarantees Input-to-State Stability (ISS) of the tracking error under Assumptions 1, 2, and 3, while utilizing the null-space to safely avoid joint limits, self-collisions, and unnatural hand configurations.

\section{Proposed Scheme}

\subsection{Singularity-Robust CLIK Control Law}
To prevent numerical instability near kinematic singularities, we employ a Damped Least-Squares (DLS) Jacobian inverse. The proposed control law is formulated as:
\begin{equation}
\begin{split}
\dot{q}(t) &= J^{\dagger}_{DLS}(q(t)) \Big[ v_{d,filt}(t) + K_p e(t) \\
&\quad + K_d e(t - \tau_{filt}(t)) \Big] \\
&\quad + \left( I_n - J^{\dagger}_{DLS}(q(t))J(q(t)) \right) \dot{q}_0(t)
\end{split}
\end{equation}
where:
\begin{itemize}
\item $J^{\dagger}_{DLS}(q) = J^T(JJ^T + \lambda^2 I_m)^{-1} \in \mathbb{R}^{n \times m}$ is the DLS pseudo-inverse, with $\lambda > 0$ being the damping factor.
\item $K_p, K_d \in \mathbb{R}^{m \times m}$ are positive-definite diagonal gain matrices.
\item $K_d e(t - \tau_{filt}(t))$ is the discrete delay-feedback compensation term.
\item $v_{d,filt}(t) = \frac{d}{dt} x_d(t - \tau_{filt}(t)) = v_d(t - \tau_{filt}(t))(1 - \dot{\tau}_{filt}(t))$ is the analytical derivative of the filtered reference trajectory.
\item $\dot{q}_0(t) \in \mathbb{R}^n$ is the joint velocity vector designed for secondary task optimization.
\end{itemize}

Substituting the control law into the differential kinematics yields the closed-loop tracking error dynamics:
\begin{equation}
\dot{e}(t) = -E(e_o) \left[ K_p e(t) + K_d e(t - \tau_{filt}(t)) \right] + d_t(t)
\end{equation}
where $E(e_o) = \text{diag}\left(I_3, \frac{1}{2}(\eta_e I_3 + \epsilon_e^\times), I_{5}\right)$ is the orientation scaling matrix, and $d_t(t)$ is the aggregate perturbation term:
\begin{equation}
d_t(t) = d_s(t) + d_l(t)
\end{equation}
with:
\begin{itemize}
\item $d_s(t)$ being the damping perturbation defined as:
\begin{equation*}
\begin{split}
d_s(t) &= \left( J J^{\dagger}_{DLS} - I_m \right) \Big[ v_{d,filt}(t) \\
&\quad + K_p e(t) + K_d e(t - \tau_{filt}(t)) \Big]
\end{split}
\end{equation*}
\item $d_l(t) = -J(q)\left( I_n - J^{\dagger}_{DLS}(q)J(q) \right)\dot{q}_0(t)$ being the workspace leakage originating from the DLS null-space projection.
\end{itemize}

To rigorously address the nonlinearities in $E(e_o)$, we define the linearization residual matrix mapping $\Delta E(e_o) = I_m - E(e_o)$. This allows the exact closed-loop error kinematics to be formulated as:
\begin{equation}
\begin{split}
\dot{e}(t) &= -K_p e(t) - K_d e(t - \tau_{filt}(t)) + d_t(t) \\
&\quad + \Delta E(e_o) \left[ K_p e(t) + K_d e(t - \tau_{filt}(t)) \right]
\end{split}
\end{equation}

\textbf{Remark 1 (Local ISS)}: Within the region of attraction $\Omega_o = \{ e_o \in \mathbb{R}^3 \mid \eta_e > 0 \}$ (corresponding to rotation errors of less than $180^\circ$), the residual mapping $\Delta E(e_o) = I_m - E(e_o)$ converges to zero near the equilibrium state. By restricting the stability analysis to a local neighborhood where $\|\Delta E(e_o)\|$ is sufficiently small, its higher-order perturbative effects on the LKF derivative are strictly dominated by the negative-definite matrix $\Omega_{ISS}$. Thus, the established stability property is Local Input-to-State Stability (Local ISS).

\subsection{Stability Analysis via Lyapunov-Krasovskii Functional}
To prove the Local ISS of the closed-loop error system, we construct the following candidate Lyapunov-Krasovskii Functional (LKF) $V(t)$:
\begin{equation}
\begin{split}
V(t) &= e^T(t) P e(t) + \int_{t - \tau_{filt}(t)}^{t} e^T(s) Q e(s) ds \\
&\quad + \tau_m \int_{-\tau_m}^{0} \int_{t + \theta}^{t} \dot{e}^T(s) R \dot{e}(s) ds d\theta
\end{split}
\end{equation}
where $P, Q, R \in \mathbb{R}^{m \times m}$ are diagonal positive-definite matrices (hence $P K_p = K_p P$, $P K_d = K_d P$).

Differentiating $V(t)$ along the trajectory of the perturbed error dynamics yields:
\begin{equation}
\begin{split}
\dot{V}(t) &= 2 e^T(t) P \left[ -K_p e(t) - K_d e(t - \tau_{filt}(t)) + d_t(t) \right] \\
&\quad + e^T(t) Q e(t) - \tau_m \int_{t - \tau_m}^{t} \dot{e}^T(s) R \dot{e}(s) ds \\
&\quad - (1 - \dot{\tau}_{filt}(t)) e^T(t - \tau_{filt}(t)) Q e(t - \tau_{filt}(t)) \\
&\quad + \tau_m^2 \dot{e}^T(t) R \dot{e}(t)
\end{split}
\end{equation}

Using Young's inequality, the bilinear perturbation term is bounded by:
\begin{equation}
2 e^T(t) P d_t(t) \le \epsilon_1 e^T(t) e(t) + \frac{1}{\epsilon_1} d_t^T(t) P^2 d_t(t)
\end{equation}
where $\epsilon_1 > 0$.

Let $\phi_e(t) = -K_p e(t) - K_d e(t - \tau_{filt}(t))$. The acceleration-related LKF term contains:
\begin{equation}
\dot{e}^T(t) R \dot{e}(t) = \left[ \phi_e(t) + d_t(t) \right]^T R \left[ \phi_e(t) + d_t(t) \right]
\end{equation}
Applying Young's inequality to the cross-term $2 \phi_e^T(t) R d_t(t)$ yields:
\begin{equation}
\dot{e}^T(t) R \dot{e}(t) \le (1 + \epsilon_2) \phi_e^T(t) R \phi_e(t) + \left( 1 + \frac{1}{\epsilon_2} \right) d_t^T(t) R d_t(t)
\end{equation}
where $\epsilon_2 > 0$.

To handle the delay-dependent integral term, we partition the integration interval into two segments: $[t - \tau_m, t - \tau_{filt}(t)]$ and $[t - \tau_{filt}(t), t]$. Applying Jensen's inequality to each segment yields:
\begin{equation}
-\tau_m \int_{t - \tau_{filt}(t)}^{t} \dot{e}^T(s) R \dot{e}(s) ds \le -\frac{\tau_m}{\tau_{filt}(t)} a^T(t) R a(t)
\end{equation}
\begin{equation}
-\tau_m \int_{t - \tau_m}^{t-\tau_{filt}(t)} \dot{e}^T(s) R \dot{e}(s) ds \le -\frac{\tau_m}{\tau_m - \tau_{filt}(t)} b^T(t) R b(t)
\end{equation}
where $a(t) = e(t) - e(t-\tau_{filt}(t))$ and $b(t) = e(t-\tau_{filt}(t)) - e(t-\tau_m)$.

Let $\eta(t) = [e^T(t), e^T(t-\tau_{filt}(t)), e^T(t-\tau_m)]^T \in \mathbb{R}^{3m}$ be the augmented state vector. Applying the \textit{reciprocal convexity lemma} (Park et al., 2011), if there exists a matrix $S \in \mathbb{R}^{m \times m}$ such that $\begin{bmatrix} R & S \\ S^T & R \end{bmatrix} \ge 0$, the partitioned segments can be bounded. We obtain:
\begin{equation}
\dot{V}(t) \le -\eta^T(t) \Omega_{ISS} \eta(t) + \gamma \|d_t(t)\|^2
\end{equation}
where the ISS-gain parameter $\gamma$ is explicitly formulated as:
\begin{equation}
\gamma = \frac{1}{\epsilon_1}\lambda_{max}(P^2) + \tau_m^2 \left( 1 + \frac{1}{\epsilon_2} \right) \lambda_{max}(R)
\end{equation}
and the $3m \times 3m$ matrix $\Omega_{ISS}$ is defined as:
\begin{equation}
\Omega_{ISS} = \Omega_1 + \Omega_2 - \Omega_3
\end{equation}
where:
\begin{align*}
\Omega_1 &= \begin{bmatrix} 2 P K_p - Q - \epsilon_1 I_m & P K_d & 0 \\ * & (1-d)Q & 0 \\ * & * & 0 \end{bmatrix}, \\
\Omega_2 &= \begin{bmatrix} R & S-R & -S \\ * & 2R-S-S^T & S-R \\ * & * & R \end{bmatrix}, \\
\Omega_3 &= \tau_m^2 (1 + \epsilon_2) \begin{bmatrix} K_p^T R K_p & K_p^T R K_d & 0 \\ * & K_d^T R K_d & 0 \\ * & * & 0 \end{bmatrix}.
\end{align*}

By solving the LMI $\Omega_{ISS} > 0$ under the constraint $\begin{bmatrix} R & S \\ S^T & R \end{bmatrix} \ge 0$, the derivative satisfies $\dot{V}(t) \le -\lambda_{min}(\Omega_{ISS}) \|\eta(t)\|^2 + \gamma \|d_t(t)\|^2$, proving the Local Input-to-State Stability (Local ISS) of the tracking error.

\subsection{Null-Space Projection for Joint Limit, Collision, and Posture Avoidance}
The secondary joint velocity vector $\dot{q}_0(t)$ is designed to push the robot joints away from physical boundaries, avoid self-collisions, and guide the redundant hand joints toward natural anthropomorphic configurations to resolve the $19 - 5 = 14$ DoFs under-constraint. We define the objective function $H(q) \in \mathbb{R}$ to be minimized:
\begin{equation}
\begin{split}
H(q) &= w_{lim} \sum_{i=1}^{n} \left( \frac{q_i - \bar{q}_i}{q_{i,max} - q_{i,min}} \right)^2 + H_{coll}(q) \\
&\quad + w_{post} \sum_{j=7}^{25} \left( q_j - q_{j,nom} \right)^2
\end{split}
\end{equation}
where:
\begin{itemize}
\item $q_{i,max}, q_{i,min}$ are the physical boundaries of the $i$-th joint, and $\bar{q}_i = \frac{q_{i,max} + q_{i,min}}{2}$.
\item $q_{j,nom}$ denotes the nominal natural posture joint angle of the $j$-th hand joint, calibrated from the operator's relaxed hand position.
\item $w_{lim}$ and $w_{post}$ are positive weight parameters.
\end{itemize}

To prevent gradient explosion and chattering in discrete-time execution (100 Hz), the collision potential $H_{coll}(q)$ is active only within a finite influence distance $d_{inf} > 0$:
\begin{equation}
H_{coll}(q) = \sum_{a < b} h_{ab}(q)
\end{equation}
where:
\begin{equation*}
h_{ab}(q) = \begin{cases}
\frac{\sigma}{2} \left( \frac{1}{d_{ab}(q)} - \frac{1}{d_{inf}} \right)^2, & \text{if } d_{ab}(q) \le d_{inf} \\
0, & \text{if } d_{ab}(q) > d_{inf}
\end{cases}
\end{equation*}
and $d_{ab}(q)$ is the distance between capsule representations of link $a$ and link $b$.

To ensure the perturbation vector $d_l(t)$ remains bounded in the LKF analysis, the gradient mapping is saturated to bound the norm of the potential vector:
\begin{equation}
\dot{q}_0(t) = \text{sat}_{v_{max}} \left( -\alpha \nabla H(q) \right)
\end{equation}
where $\text{sat}_{v_{max}}(\cdot)$ is an element-wise saturation function bounding each component to $\pm v_{max}$. Consequently, the Euclidean norm of the secondary velocity is strictly bounded ($\| \dot{q}_0(t) \| \le \sqrt{n} v_{max}$). This guarantees that the DLS leakage perturbation $d_l(t)$ remains bounded globally, preserving the ISS conditions.

\textbf{Remark 2 (SVD-based Null-Space Leakage)}: Applying singular value decomposition (SVD) to the Jacobian $J = U \Sigma V^T$, the workspace leakage introduced by the DLS projection is given by $J \left( I_n - J^{\dagger}_{DLS}J \right) = U \text{diag}\left( \frac{\sigma_i \lambda^2}{\sigma_i^2 + \lambda^2} \right) V^T$. By AM-GM inequality, each diagonal entry is bounded by $\frac{\lambda}{2}$ for all $\sigma_i$. As the system approaches a deep singularity ($\sigma_i \to 0$), the leakage converges to zero. The maximum leakage is bounded globally by $\mathcal{O}(\lambda)$ at near-singular configurations ($\sigma_i \approx \lambda$), ensuring that the secondary tasks do not corrupt workspace accuracy.

\subsection{Practical Implementation Details}
In practice, the delay $\tau(t)$ is calculated in real-time using network packet timestamps: $\tau(t) = t_{recv} - t_{send}$. Because raw latency measurements contain high-frequency noise (jitter), the derivative $\dot{\tau}(t)$ cannot be computed directly via numerical differentiation. 

To resolve this issue, the controller implements a continuous-time Low-Pass Filter (LPF):
\begin{equation}
\tau_{filt}(s) = \frac{1}{T_f s + 1} \tau(s)
\end{equation}
To strictly guarantee Assumption 2 ($|\dot{\tau}_{filt}(t)| \le d$), the filter time constant $T_f$ cannot be chosen arbitrarily. Given the LPF time-domain dynamics $\dot{\tau}_{filt} = \frac{1}{T_f} (\tau - \tau_{filt})$ and the bounded jitter $\|n(t)\| \le n_{max}$ from Assumption 3, $T_f$ must be selected such that:
\begin{equation}
T_f \ge \frac{\max |\tau(t) - \tau_{filt}(t)|}{d}
\end{equation}
In practical implementation, $T_f$ is conservatively tuned above this theoretical lower bound to ensure smooth command velocities.

Furthermore, computing the collision gradient $\nabla H_{coll}(q)$ requires the real-time calculation of the closest-point Jacobians $J_{ab}(q)$ for all active capsule pairs:
\begin{equation}
\frac{\partial d_{ab}(q)}{\partial q} = \hat{n}_{ab}^T J_{ab}(q)
\end{equation}
where $\hat{n}_{ab}$ is the unit normal vector between the closest points. To ensure execution at 100 Hz, spatial hashing or bounding-volume hierarchy (BVH) algorithms provided by the Isaac Sim API are utilized to quickly cull inactive collision pairs ($d_{ab} > d_{inf}$), isolating the Jacobian computations only to the active proximity sets. This CLIK controller operates as a high-frequency (100 Hz) teleoperation bridge in the Isaac Sim simulator, allowing the human operator to generate stable, chatter-free demonstration data for downstream Diffusion Policy training.

\begin{algorithm}[H]
\caption{Singularity-Robust and Delay-Robust Unified CLIK Control}
\label{alg:clik}
\begin{algorithmic}[1]
\REQUIRE Current joint configuration $q(t) \in \mathbb{R}^{25}$, raw communication delay $\tau(t)$, desired workspace trajectory $x_d(t) \in \mathcal{X}$, nominal synergy hand posture $q_{nom} \in \mathbb{R}^{19}$, LPF time constant $T_f$, gains $K_p, K_d, \alpha, \sigma$, damping $\lambda$.
\ENSURE Command joint velocity vector $\dot{q}(t) \in \mathbb{R}^{25}$.
\STATE Measure current packet latency $\tau(t) = t_{recv} - t_{send}$.
\STATE Update filtered delay state $\dot{\tau}_{filt}(t) = \frac{1}{T_f} (\tau(t) - \tau_{filt}(t))$.
\STATE Compute filtered reference trajectory $x_{d,filt}(t) = x_d(t - \tau_{filt}(t))$.
\STATE Compute reference velocity $v_{d,filt}(t) = v_d(t - \tau_{filt}(t))(1 - \dot{\tau}_{filt}(t))$.
\STATE Compute tracking error $e(t) = [e_p^T(t), e_o^T(t), e_f^T(t)]^T$ using unit quaternion for orientation.
\STATE Compute analytical system Jacobian $J(q(t)) \in \mathbb{R}^{11 \times 25}$.
\STATE Calculate DLS pseudo-inverse $J^{\dagger}_{DLS} = J^T\left(J J^T + \lambda^2 I_{11}\right)^{-1}$.
\STATE Query spatial culling interface (BVH/Hashing) for active capsule pairs ($d_{ab} \le d_{inf}$).
\STATE Compute secondary gradient $\nabla H = \nabla H_{lim} + \nabla H_{coll} + \nabla H_{post}$.
\STATE Compute saturated null-space velocity $\dot{q}_0(t) = \text{sat}_{v_{max}} \left( -\alpha \nabla H \right)$.
\STATE Compute task space control input $u_e(t) = v_{d,filt} + K_p e(t) + K_d e(t - \tau_{filt})$.
\STATE Calculate command joint velocity vector:
$$\dot{q}(t) = J^{\dagger}_{DLS} u_e(t) + \left( I_{25} - J^{\dagger}_{DLS} J \right) \dot{q}_0(t)$$
\RETURN $\dot{q}(t)$
\end{algorithmic}
\end{algorithm}

\begin{thebibliography}{33}
\bibitem{ref_diffpolicy3d}
Y. Ze, G. Zhang, K. Zhang, C. Hu, M. Wang, and H. Xu, ``3d diffusion policy: Generalizable visuomotor policy learning via simple 3d representations,'' \textit{arXiv preprint arXiv:2403.03954}, 2024.
\bibitem{ref_flowdiff}
S. Noh, D. Nam, K. Kim, G. Lee, Y. Yu, R. Kang, and K. Lee, ``3d flow diffusion policy: Visuomotor policy learning via generating flow in 3d space,'' \textit{arXiv preprint arXiv:2509.18676}, 2025.
\bibitem{ref_egocentric3d}
A. Prakash, R. Tu, M. Chang, and S. Gupta, ``3d hand pose estimation in everyday egocentric images,'' in \textit{Proc. Eur. Conf. Comput. Vis. (ECCV)}, 2024, pp. 183--202.
\bibitem{ref_cybernetics_armhand}
S. Li, N. Hendrich, H. Liang, P. Ruppel, C. Zhang, and J. Zhang, ``A dexterous hand-arm teleoperation system based on hand pose estimation and active vision,'' \textit{IEEE Trans. Cybernetics}, vol. 54, no. 3, pp. 1417--1428, Mar. 2022.
\bibitem{ref_a2j_trans}
C. Jiang, Y. Xiao, C. Wu, M. Zhang, J. Zheng, Z. Cao, and J. T. Zhou, ``A2j-transformer: Anchor-to-joint transformer network for 3d interacting hand pose estimation from a single rgb image,'' in \textit{Proc. IEEE/CVF Conf. Comput. Vis. Pattern Recognit. (CVPR)}, 2023, pp. 8846--8855.
\bibitem{ref_amp}
X. B. Peng, Z. Ma, P. Abbeel, S. Levine, and A. Kanazawa, ``Amp: Adversarial motion priors for stylized physics-based character control,'' \textit{ACM Trans. Graphics (ToG)}, vol. 40, no. 4, pp. 1--20, 2021.
\bibitem{ref_xin_retarget}
C. Xin, M. Yu, Y. Jiang, Z. Zhang, and X. Li, ``Analyzing Key Objectives in Human-to-Robot Retargeting for Dexterous Manipulation,'' \textit{IEEE Robotics and Automation Practice}, 2026.
\bibitem{ref_anyteleop}
Y. Qin, W. Yang, B. Huang, K. Van Wyk, H. Su, X. Wang, Y. Chao, D. Fox, and X. Wang, ``Anyteleop: A general vision-based dexterous robot arm-hand teleoperation system,'' \textit{arXiv preprint arXiv:2307.04577}, 2023.
\bibitem{ref_egohand_bench}
Z. Fan, T. Ohkawa, L. Yang, N. Lin, Z. Zhou, S. Zhou, and A. Yao, ``Benchmarks and challenges in pose estimation for egocentric hand interactions with objects,'' in \textit{Proc. Eur. Conf. Comput. Vis. (ECCV)}, 2024, pp. 428--448.
\bibitem{ref_zero_shot_dex}
Z. Zhao, Z. Li, Y. Ou, and M. Qi, ``Closing the reality gap: Zero-shot sim-to-real deployment for dexterous force-based grasping and manipulation,'' \textit{arXiv preprint arXiv:2607.04940}, 2026.
\bibitem{ref_deformer}
Q. Fu, X. Word, R. Xu, J. C. Niebles, and K. M. Kitani, ``Deformer: Dynamic fusion transformer for robust hand pose estimation,'' in \textit{Proc. IEEE/CVF Int. Conf. Comput. Vis. (ICCV)}, 2023, pp. 23600--23611.
\bibitem{ref_deltadorsal}
W. Huang, S. Pei, L. Zou, E. J. Gonzalez, I. Chatterjee, and Y. Zhang, ``DeltaDorsal: Enhancing Hand Pose Estimation with Dorsal Features in Egocentric Views,'' in \textit{Proc. CHI Conf. Human Factors Comput. Syst.}, 2026, pp. 1--16.
\bibitem{ref_dexhil}
Y. Han, Z. Chen, Y. Zhao, C. Xu, Y. Shao, Y. Peng, and W. Lian, ``DexHiL: A Human-in-the-Loop Framework for Vision-Language-Action Model Post-Training in Dexterous Manipulation,'' \textit{arXiv preprint arXiv:2603.09121}, 2026.
\bibitem{ref_dexsim2real}
Z. Zeng, F. Ding, H. Yang, X. Li, and Y. Liao, ``DexSim2Real: Foundation Model-Guided Sim-to-Real Transfer for Generalizable Dexterous Manipulation,'' \textit{arXiv preprint arXiv:2605.05241}, 2026.
\bibitem{ref_tase_tele}
B. Jing, J. Na, H. Duan, X. Wang, and Y. Huang, ``A USDE-based control for bilateral teleoperation systems with time-varying delays,'' \textit{IEEE Trans. Autom. Sci. Eng.}, vol. 22, pp. 407--417, 2024.
\bibitem{ref_tcyb_nn}
S. Zhang, S. Yuan, X. Yu, L. Kong, Q. Li, and G. Li, ``Adaptive neural network fixed-time control design for bilateral teleoperation with time delay,'' \textit{IEEE Trans. Cybernetics}, vol. 52, no. 9, pp. 9756--9769, Sep. 2021.
\bibitem{ref_tmech_manip}
S. Zheng, Y. Zha, C. K. Ahn, S. Lu, and B. Song, ``Constrained finite-time output regulation for robot manipulators with control input delay,'' \textit{IEEE/ASME Trans. Mechatronics}, vol. 30, no. 6, pp. 6654--6666, Dec. 2025.
\bibitem{ref_tcst_delay}
A. Bertino, P. Naseradinmousavi, and M. Krstic, ``Delay-adaptive control of a 7-DOF robot manipulator: Design and experiments,'' \textit{IEEE Trans. Control Syst. Technol.}, vol. 30, no. 6, pp. 2506--2521, Nov. 2022.
\bibitem{ref_access_dual}
F. A. Chicaiza, E. Slawinski, and V. Mut, ``Dual-arm mobile manipulator teleoperation with coupled rate-position mapping under time-varying delays,'' \textit{IEEE Access}, 2025.
\bibitem{ref_neuro_imped}
Z. Xu, X. Li, S. Li, H. Wu, and X. Zhou, ``Dynamic neural networks based adaptive optimal impedance control for redundant manipulators under physical constraints,'' \textit{Neurocomputing}, vol. 471, pp. 149--160, 2022.
\bibitem{ref_tcyb_event}
Y. Li, C. Li, K. Liu, J. Dong, and R. Johansson, ``Event-Triggered Control and Communication for Single-Master Multislave Teleoperation Systems With Try-Once-Discard Protocol,'' \textit{IEEE Trans. Cybernetics}, 2025.
\bibitem{ref_cac_event}
C. Li, Y. Li, J. Dong, and H. Wang, ``Event-triggered control of teleoperation systems with time-varying delays,'' in \textit{Proc. China Automation Congress (CAC)}, 2021, pp. 1537--1542.
\bibitem{ref_tro_grasp}
C. Guo, X. Chen, Z. Zeng, Z. Guo, Y. Li, H. Xiao, and H. Lu, ``Grasp like humans: Learning generalizable multi-fingered grasping from human proprioceptive sensorimotor integration,'' \textit{IEEE Trans. Robotics}, 2025.
\bibitem{ref_ojies_coor}
E. Slawinski, F. A. Chicaiza, F. G. Rossomando, V. Mut, and J. Moreno-Valenzuela, ``Lyapunov-stable Neural Network Control for Position and Force Coordination in Delayed Bilateral Teleoperation Systems,'' \textit{IEEE Open J. Ind. Electron. Soc.}, 2026.
\bibitem{ref_tsmc_nn}
W. M. Wang and Y. W. Wang, ``New Criteria on Stability Analysis of Generalized Neural Networks With Time-Varying Delays,'' \textit{IEEE Trans. Syst., Man, Cybern.: Systems}, 2025.
\bibitem{ref_tro_rcm}
T. Kastritsi, T. P. Semetzidis, and Z. Doulgeri, ``Passive bilateral surgical teleoperation with RCM and spatial constraints in the presence of time delays,'' \textit{IEEE Trans. Robotics}, vol. 41, pp. 612--627, 2024.
\bibitem{ref_ijsr_tele}
C. Wang, R. Luo, J. P. Whitney, and T. Padir, ``Practical approaches towards transparent and stable bilateral teleoperation under time-varying network delay,'' \textit{Int. J. Soc. Robotics}, vol. 17, no. 3, pp. 409--426, 2025.
\bibitem{ref_access_pred}
L. Chan, Y. Liu, Q. Huang, and P. Wang, ``Robust adaptive observer-based predictive control for a non-linear delayed bilateral teleoperation system,'' \textit{IEEE Access}, vol. 10, pp. 52294--52305, 2022.
\bibitem{ref_iecon_sliding}
S. Islam and A. Sunda-Meya, ``Robust sliding mode based finite-time bilateral shared teleoperation system with unsymmetrical time-varying delay,'' in \textit{Proc. IECON 48th Annu. Conf. IEEE Ind. Electron. Soc.}, 2022, pp. 1--6.
\bibitem{ref_tcs_haptic}
Y. Liu, D. Xiong, L. Wang, and C. K. Zhang, ``Stability analysis of haptic systems with time-varying delay via a delay-product-type Lyapunov–Krasovskii functional,'' \textit{IEEE Trans. Circuits Syst. II: Express Briefs}, vol. 69, no. 11, pp. 4339--4343, Nov. 2022.
\bibitem{ref_tsmc_linear}
C. R. Wang, Y. He, K. Z. Liu, and C. K. Zhang, ``Stability analysis of linear systems with a time-varying delay via less conservative methods,'' \textit{IEEE Trans. Syst., Man, Cybern.: Systems}, vol. 55, no. 3, pp. 1972--1983, 2024.
\bibitem{ref_tcyb_piece}
B. Kaviarasan, O. M. Kwon, M. J. Park, and R. Sakthivel, ``Stabilization of periodic piecewise time-varying systems with time-varying delay under multiple cyber attacks: An augmented Lyapunov functional approach,'' \textit{IEEE Trans. Cybernetics}, vol. 54, no. 8, pp. 4389--4401, Aug. 2023.
\bibitem{ref_tcst_delayed_force}
A. Gutierrez-Giles, A. Rodriguez-Angeles, P. Schleer, and M. A. Arteaga, ``Time-varying delayed bilateral teleoperation without force/velocity measurements,'' \textit{IEEE Trans. Control Syst. Technol.}, vol. 32, no. 3, pp. 780--792, May 2023.
\end{thebibliography}

\end{document}
